\section{结论}
\label{sec:conclusion}

本文面向障碍物、柱状禁飞区、高度边界和风险热点共同作用的强约束三维低空路径规划问题，提出 CVF-AE。其核心不在于以额外惩罚项替代原有目标，而是将局部约束信息前移至候选生成阶段，并保持评价函数和 Deb 可行性优先准则不变。具体而言，约束感知初始化先将部分种群置于潜在可行走廊附近；状态调度再依据种群可行性、违反程度和搜索进展决定介入时机；CVF 将障碍物、禁飞区、风险、高度、边界和转弯信息映射为控制点修正方向，并通过逐维与整体步长截断限制其影响范围；最后，CVF 候选须同时满足局部 Deb 更优和综合适应度不增才可替代内部基础候选。由此，算法形成了“引导进入--按需介入--受限修正--保守保留”的闭环，使约束信息能够改善候选质量，而不演变为对全种群的强制修复。

三类约束逐步收紧的低空场景支持该闭环设计。在本文固定离散采样判定下，CVF-AE 在所选代表性基线中三场景可行率均为 $1.00$，以 $1.33$ 的平均排名取得整体最优，并在约束最密集的 Scene 3 中保持较优的最终适应度。消融结果进一步说明各环节的必要性：移除完整候选增强闭环后，三个场景的可行率降至 $0.63$、$0$ 和 $0$；去除约束感知初始化时，Scene 3 的可行率降至 $0.47$ 且 CVF 触发数显著增加；去除状态调度或稀疏可行性保持也会损害强约束场景中的路径质量。完整方法相对于同配置 w/o-CVF 版本的三场景平均额外评价次数约为 $0.44\%$，表明方法收益来自有针对性的稀疏介入，而非全种群额外修复。

本文结论限于静态合成场景及固定采样下的数值可行性。后续可在真实地图与动态障碍环境中结合连续碰撞验证、在线重规划和飞行动力学/轨迹跟踪控制，进一步检验 CVF-AE 的工程适用性。
